% STATUT : À RÉDIGER, résultats produits (benchmark copule, benchmark
% Formule Standard). Figure W_benchmark_sf.png, S18_benchmark_copule.png.
\chapter{Ce que les cadres existants expriment de l'effet DORA}
\label{ch:benchmarks}

\textit{[Chapitre à rédiger.]}

% MESSAGE CENTRAL, en une ligne de lecture :
%   quelle FRACTION de l'effet DORA (l'écart de capital entre un profil
%   conforme et un profil non conforme) chaque cadre est capable
%   d'exprimer ?
%     Formule Standard ......  0 %   (aucune sensibilité au niveau de maîtrise)
%     Copule ................ 62 %   (dépendance symétrique, sans direction)
%     Cascade dirigée ...... 100 %   (par construction, c'est la référence)
%
% PLAN :
%  1. Formule Standard : SCR_op = min(0,3 x BSCR ; 0,03 x max(primes,
%     provisions)). Montrer qu'elle est structurellement aveugle à DORA.
%  2. Copule : ce qu'elle capture (dépendance de queue) et ce qu'elle rate
%     (la direction). Sa priorité de remédiation suit les RÉCEPTACLES,
%     la cascade suit les SOURCES : c'est la différence opérationnelle.
%  3. NOTE DE PÉRIMÈTRE : la copule apparaît ici comme POINT DE
%     COMPARAISON uniquement. Le modèle à briques (LDA 4 briques + Gumbel
%     + latente) est abandonné et n'est plus un modèle du mémoire.
%  4. Discussion : ce benchmark est ce qui rend l'apport lisible pour un
%     praticien, il doit être mis en avant, pas relégué en annexe.
